\documentclass{article}
\usepackage{amsmath}
\usepackage{amsfonts}
\usepackage{amssymb}
\usepackage[top=1.75cm,bottom=1.5cm,right=2cm,left=2cm]{geometry}
\setlength{\headheight}{22pt}
\usepackage{enumitem}
\usepackage{fancyhdr} % For customizing headers and footers

%---------------- Header
\pagestyle{fancy}
\fancyhf{} % Clear default header and footer
\fancyhead[L]{{\footnotesize CUNEF Universidad \\ Bachelor in Philosophy, Politics and Economics}} 
\fancyhead[R]{{\footnotesize A. G. Azze \\ \date{2024}}} 
\setlength{\headsep}{1.2cm}
\fancypagestyle{plain}{
  \fancyhf{}
  \fancyhead[L]{{\footnotesize CUNEF Universidad \\ Bachelor in Philosophy, Politics and Economics}} 
  \fancyhead[R]{{\footnotesize A. G. Azze \\ \date{2025}}} 
}

%---------------- Title
\title{Math Essentials \\ Unit 4 - Logarithms and Exponentials}
\author{}
\date{}

\begin{document}
\maketitle

\section*{Exercises}

\begin{enumerate}[label=\textbf{\arabic*.}, leftmargin=*, itemsep=0.6em]

\item In chemistry, \(\mathrm{pH}\) is a measure of acidity and is given by \(\mathrm{pH}=-\log(H^+)\), where \(H^+\) is the hydrogen ion concentration (in moles of hydrogen per liter of solution). Determine \(H^+\) if \(\mathrm{pH}=4\).

\item The decibel level to measure volume is given by
\[
D \;=\; 10\,\log_{10}\!\left(\frac{I}{I_0}\right), \qquad I_0 = 10^{-12}\ \mathrm{W/m^2}.
\]
\begin{enumerate}[label=(\alph*)]
  \item An alarm measures \(D=120\) dB. Find its intensity \(I\).
  \item A jet has intensity \(I=8\cdot 3\cdot 10^{2}\ \mathrm{W/m^2}\). Find its decibel level \(D\).
\end{enumerate}

\item The population \(P\) of a small town is modeled by \(P=1650e^{0.5t}\) where \(t\) is measured in years. In approximately how many years will the town’s population reach \(20{,}000\)?

\item Atmospheric pressure \(P\) in pounds per square inch is represented by \(P=14.7e^{-0.21x}\), where \(x\) is the number of miles above sea level. To the nearest foot, how high is the peak of a mountain with an atmospheric pressure of \(8.369\) psi? (Hint: \(5280\) feet \(=1\) mile.)

\item The magnitude \(M\) of an earthquake is represented by \(M=\dfrac{2}{3}\log\!\left(\dfrac{E}{E_0}\right)\) where \(E\) is the energy (in joules) released and \(E_0=10^{4.4}\) is the assigned minimal measure released by an earthquake. To the nearest hundredth, what is the magnitude of an earthquake releasing \(1.4\cdot 10^{13}\) joules?

\item What is the value of  \(b^{\log_b x}\)?

\item Solve  \(y=Ae^{kt}\) for $t$.

\item Recall the compound interest formula \(A=a\big(1+\tfrac{r}{k}\big)^{kt}\). Use the definition of a logarithm along with properties of logarithms to solve the formula for time \(t\).

\item Newton’s Law of Cooling states that the temperature \(T\) of an object at any time \(t\) can be described by
\(T=T_s+(T_0-T_s)e^{-kt}\),
where \(T_s\) is the surrounding temperature, \(T_0\) is the object’s initial temperature, and \(k\) is the cooling rate. Use logarithms to solve the formula for time \(t\) so that \(t\) is a single logarithm.

\end{enumerate}

\newpage
\section*{Solutions}

\begin{enumerate}[label=\textbf{\arabic*.}, leftmargin=*, itemsep=0.45em]

\item \(\mathrm{pH}=4\Rightarrow -\log(H^+)=4\Rightarrow \log(H^+)=-4\Rightarrow H^+=10^{-4}\) mol/L.

\item 
\begin{enumerate}[label=(\alph*)]
  \item From \(120 = 10\log_{10}\!\left(\dfrac{I}{10^{-12}}\right)\) we get
  \[
  \log_{10}\!\left(\frac{I}{10^{-12}}\right)=12
  \;\Longrightarrow\;
  \frac{I}{10^{-12}}=10^{12}
  \;\Longrightarrow\;
  I = 10^{-12}\cdot 10^{12}=1\ \mathrm{W/m^2}.
  \]
  \item Using \(D=10\log_{10}(I/I_0)\) with \(I=8\cdot 3\cdot 10^{2}\):
  \[
  D
  =10\log_{10}(24\cdot 10^{14})
  =10\big(\log_{10}(24)+\log_{10}(10^{14})\big)
  = 10\big(\log_{10}(24) + 14\big)\ \text{dB}.
  \]
\end{enumerate}

\item \(1650e^{0.5t}=20000\Rightarrow e^{0.5t}=\dfrac{20000}{1650}\Rightarrow
0.5t=\ln\!\left(\dfrac{20000}{1650}\right)\Rightarrow t\approx 4.99\ \text{years}.\)

\item \(8.369=14.7e^{-0.21x}\Rightarrow e^{-0.21x}=\dfrac{8.369}{14.7}\Rightarrow
x=-\dfrac{1}{0.21}\ln\!\left(\dfrac{8.369}{14.7}\right)\approx 2.6824\ \text{miles}.\)
Height \(\approx 2.6824\times 5280\approx 14{,}163\ \text{ft}.\)

\item \(M=\dfrac{2}{3}\log\!\left(\dfrac{E}{E_0}\right)
=\dfrac{2}{3}\log\!\left(\dfrac{1.4\cdot 10^{13}}{10^{4.4}}\right)
=\dfrac{2}{3}\left(\log\!\left(\dfrac{1.4\cdot 10^{8.613}}{10^{4.4}}\right)\right)
= \dfrac{2}{3}\left(\log\!\left(1.4\right) + 10^{8.6}\right)
\approx \dfrac{2}{3}\cdot 8.7461\approx 5.83.\)

\item Let \(y=\log_b x\). By definition, \(b^y=x\). Hence \(b^{\log_b x}=x\).

\item From \(y=Ae^{kt}\): \(\dfrac{y}{A}=e^{kt}\Rightarrow \ln\!\left(\dfrac{y}{A}\right)=kt\Rightarrow
t=\dfrac{1}{k}\ln\!\left(\dfrac{y}{A}\right).\)

\item \(A=a\!\left(1+\dfrac{r}{k}\right)^{kt}\Rightarrow \dfrac{A}{a}=\!\left(1+\dfrac{r}{k}\right)^{kt}\Rightarrow
\log\!\left(\dfrac{A}{a}\right)=kt\,\log\!\left(1+\dfrac{r}{k}\right)\Rightarrow
t=\dfrac{\log(A/a)}{k\,\log(1+r/k)}\) (any log base).

\item \(T=T_s+(T_0-T_s)e^{-kt}\Rightarrow \dfrac{T-T_s}{T_0-T_s}=e^{-kt}\Rightarrow
\ln\!\left(\dfrac{T-T_s}{T_0-T_s}\right)=-kt\Rightarrow
t=-\dfrac{1}{k}\,\ln\!\left(\dfrac{T-T_s}{T_0-T_s}\right).\)

\end{enumerate}

\end{document}